\documentclass[10pt,a4paper]{amsart}
\usepackage[margin=19mm]{geometry}
\usepackage{amsmath,amssymb,mathtools}
\usepackage[scr=rsfs,cal=euler]{mathalfa}
\usepackage{xfrac}
\usepackage[unicode,hidelinks,bookmarks=false]{hyperref}
\newcommand{\ie}{i.e.\ }
\newcommand{\cf}{cf.\ }
\newcommand{\resp}{respectively\ }
\newcommand{\Subsection}{\subsection}
\providecommand{\nobreakdash}{\nobreak\hbox{-}}
\setlength{\emergencystretch}{2em}
\multlinegap0pt
\usepackage{siunitx}
\usepackage{siunitx}
\usepackage{siunitx}
\usepackage{siunitx}
\usepackage{amssymb}
\usepackage{mathtools}
\usepackage{tikz}
\usepackage{algorithm}\usepackage{algpseudocode}
\usepackage{siunitx}
\newtheorem{theorem}{Theorem}
\newcommand{\g}{\mathfrak{g}}
\newcommand{\grad}{\normalfont{\text{grad}}}
\newcommand{\h}{\mathfrak{d}}
\title{Virtual Holonomic and Nonholonomic Constraints on Lie groups}
\author{A. Anahory Simoes, A. Bloch, L. Colombo, E. Stratoglou}
\date{Source: arXiv:2312.17531v2 (CC BY 4.0)}
\begin{document}
\maketitle

\noindent\emph{Source and licence.} Virtual Holonomic and Nonholonomic Constraints on Lie groups. Version arXiv:2312.17531v2 (CC BY 4.0), distributed by arXiv under the Creative Commons Attribution 4.0 International licence, http://creativecommons.org/licenses/by/4.0/, as recorded in the arXiv licence metadata for this exact version and shown on the abstract page https://arxiv.org/abs/2312.17531v2. Full licence notice: \texttt{licenses/arxiv-2312.17531-CC-BY-4.0.txt}.\par\smallskip

\noindent\textit{Editorial scope.} Counted unit: the article's theorem nh-orthogonal:input:thm (source0.tex lines 1124--1148). The article's complete original proof of it is delivered with it (source0.tex lines 1150--1221, one \texttt{proof} environment of 72 lines). No mathematics is written, repaired or re-derived: the statement and the proof are the article's own lines, in the article's own notation, and no retained line is substituted or re-flowed.  Retained uncounted as the source-local context the proof needs: lines 1087--1093.  Nothing was omitted from any retained span, so the package carries no omission record.  No retained line contains a citation, so the wrapper carries no bibliography and no bibliography entry.  No retained line contains a figure environment, an image-inclusion command or drawing code, so no image is delivered and the package declares no asset dependency.  Editorial formatting only, in addition: an AMS \texttt{amsart} preamble that restates the article's own declarations and macros used by the retained lines, namely the environments \texttt{theorem} on separate counters, together with the standard packages those lines need.  The retained environments print in this document as Theorem 1 in order of appearance, whereas the article prints the same items with the numbering of its own full document.  The complete native source of the article as distributed in the arXiv e-print for this exact version is bundled unchanged as \texttt{sources/152233/source0.tex}.
\begin{equation}\label{nh-control:system}
    \dot{g} = T_{e}L_{g}(\xi),\,
    \dot{\xi} 
    + \nabla^{\mathfrak{g}}_{\xi}\xi 
    + \grad_{b} V(g)
    = u^{a}f_{a}.
\end{equation}

\begin{theorem}\label{nh-orthogonal:input:thm}
Assume that the control input subspace coincides with the orthogonal complement, $\mathfrak{f} = \h^\perp = \mathrm{span}\{f_1,\dots,f_m\}$, and that $\{f_a\}$ is a basis of $\h^\perp$.

Then, for any initial condition with $\xi(0)\in\h$, there exists a unique feedback 
control law $u^\star(\xi,g)$ that renders $\h$ invariant under the closed--loop 
dynamics of \eqref{nh-control:system}.  

Moreover, this feedback is given implicitly by
\begin{equation}\label{eq:nh-feedback-implicit}
    u^{\star}_a f_a 
    = \mathfrak{Q}\Big(
        \nabla^{\g}_{\xi}\xi 
        + \grad_{b} V(g)
      \Big)\in\h^\perp,
\end{equation}
and the reduced closed--loop dynamics on $\h$ are
\begin{equation}\label{eq:nh-reduced-constraint}
    \dot{\xi} 
    + \nabla^{\h}_{\xi}\xi
    + \mathfrak{P}\Big(
        \grad_{b} V(g)
      \Big) = 0,
    \qquad \xi(t)\in\h.
\end{equation}
\end{theorem}

\begin{proof}
Rewrite \eqref{nh-control:system} as
\[
\dot{\xi} 
= -\,\nabla^{\g}_{\xi}\xi 
  - \grad_{b} V(g)
  + u^{a}f_a.
\]
We seek a feedback such that, for $\xi\in\h$, the reduced dynamics on $\h$ satisfy
\[
\dot{\xi} + \nabla^\h_{\xi}\xi
+ \mathfrak{P}\Big(
    \grad_{b} V(g)
  \Big) = 0.
\]
Using $\nabla^\h_{\xi}\xi = \mathfrak{P}(\nabla^\g_{\xi}\xi)$, this is equivalent to
\[
\dot{\xi} 
+ \mathfrak{P}\Big(
    \nabla^\g_{\xi}\xi 
    + \grad_{b} V(g)
  \Big) = 0,
\qquad \xi\in\h.
\]
Substituting the controlled dynamics for $\dot\xi$ gives
\begin{align}
&-\nabla^{\g}_{\xi}\xi 
 - \grad_{b} V(g)
 + u^{a}f_a +\, \mathfrak{P}\Big(
    \nabla^\g_{\xi}\xi 
    + \grad_{b} V(g)
  \Big) = 0,
\end{align}
hence
\begin{align}
u^{a}f_a
&= \big(I-\mathfrak{P}\big)\Big(
    \nabla^{\g}_{\xi}\xi 
    + \grad_{b} V(g)
  \Big) \nonumber\\
&= \mathfrak{Q}\Big(
    \nabla^{\g}_{\xi}\xi 
    + \grad_{b} V(g)
  \Big).
\end{align}

The right--hand side lies in $\h^\perp$, and since $\{f_a\}$ is a basis of $\h^\perp$, 
this uniquely determines the coefficients $u^{\star}_a$, proving \eqref{eq:nh-feedback-implicit} 
and the existence and uniqueness of the feedback.

Substituting $u^\star$ back into the controlled dynamics yields
\[
\dot{\xi} 
+ \nabla^{\g}_{\xi}\xi 
+ \grad_{b} V(g)
= \mathfrak{Q}\Big(
    \nabla^{\g}_{\xi}\xi 
    + \grad_{b} V(g)
  \Big),
\]
that is,
\[
\dot{\xi} 
+ \mathfrak{P}\Big(
    \nabla^{\g}_{\xi}\xi 
    + \grad_{b} V(g)
  \Big) = 0.
\]
For $\xi\in\h$, this is exactly \eqref{eq:nh-reduced-constraint}.
Since the right--hand side belongs to $\h$, solutions starting in $\h$ remain 
in $\h$, so $\h$ is invariant under the closed--loop flow.
\end{proof}

\end{document}
